% Stratégie de découpage parent-enfant
\begin{tikzpicture}[node distance=6mm and 10mm]

  \node[btpbox, text width=9.6cm] (parent)
    {\textbf{Extrait parent} — 800 à 1500 unités lexicales\\[-1pt]
     \btplabel{Conservé en base relationnelle, transmis au modèle comme contexte}};

  \node[btpboxalt, text width=2.7cm, below=12mm of parent.south west, anchor=north west] (c1)
    {Enfant 1\\[-1pt]\btplabel{250--400 unités}};
  \node[btpboxalt, text width=2.7cm, right=of c1] (c2)
    {Enfant 2\\[-1pt]\btplabel{250--400 unités}};
  \node[btpboxalt, text width=2.7cm, right=of c2] (c3)
    {Enfant 3\\[-1pt]\btplabel{250--400 unités}};

  \foreach \c in {c1, c2, c3} {
    \draw[btparrowlight] (parent.south) -- (\c.north);
  }

  \node[btplabel, below=5mm of c2, text width=10cm]
    {Les enfants seuls sont indexés dans l'espace vectoriel~: ils rendent la
     recherche précise. Les parents correspondants sont ensuite récupérés,
     dédupliqués, puis envoyés au modèle~: ils fournissent un contexte
     suffisant pour rédiger.};

  \node[btpboxstrong, text width=2.4cm, right=14mm of parent] (idx) {Recherche\\sur les enfants};
  \draw[btparrow] (idx.south) |- ($(c3.east)+(0.35,0)$) -- (c3.east);

\end{tikzpicture}
